\section{历史发展}\label{sec:history}

\subsection{互反律的谱系}

朗兰兹纲领的历史远比1967年那封信悠久——它是一部\textbf{互反律推广史}：

\begin{itemize}
    \item \textbf{高斯二次互反律}（1796）：$\left(\frac{p}{q}\right)\left(\frac{q}{p}\right) = (-1)^{\frac{p-1}{2}\cdot\frac{q-1}{2}}$，控制二次域的素数分解；
    \item \textbf{阿廷互反律}（1927）：阿贝尔扩张的伽罗瓦群与广义类群的典范同构，完成了阿贝尔情形的类域论\cite{artin1927}；阿廷同时把赫克 $L$ 函数推广到非阿贝尔伽罗瓦表示，并提出\textbf{阿廷猜想}——不可约非常数表示的 $L$ 函数整性。非阿贝尔互反律自此悬而未决，成为纲领的第一推动力。
\end{itemize}

赫克在1920年代用模形式构造的 $L$ 函数、以及阿廷 $L$ 函数之间的表面相似性，构成了“两个世界由同一本账簿（$L$ 函数）记账”的最初暗示。

\subsection{自守形式的算术化（1940s--1960s）}

纲领的第二个前提是\textbf{自守形式的表示论化}：

\begin{itemize}
    \item Chevalley 引入 Idele 类群，韦伊与泰特将其发展成分析框架：泰特1950年的博士论文\cite{tate1967}用局部--整体调和分析重证赫克 $L$ 函数的函数方程，是 adèle 语言的第一座丰碑；
    \item Harish-Chandra 建立实约化群的表示论；Gelfand 学派提出“自守表示”概念，指出尖点自守形式应对应不可约表示。
\end{itemize}

到1960年代中期，几何（模形式）、算术（狄利克雷级数）与表示（群表示）三种语言已经就位，只待一个把它们焊接在一起的猜想。

\subsection{1967年的信件与纲领的诞生}

1967年1月，普林斯顿的年轻教授朗兰兹给刚到访的韦伊写下一封十七页手写信\cite{langlands1967}，提出：

\begin{enumerate}[label=(\arabic*)]
    \item \textbf{非阿贝尔类域论}（互反律）的雏形：$\operatorname{GL}_2$ 上的尖点自守表示应对应二维伽罗瓦表示，其 $L$ 函数相等；
    \item \textbf{函子性}的一般原理：对偶群 $^LG$ 之间的同态应诱导自守表示之间的转移。
\end{enumerate}

朗兰兹随后在《Euler Products》\cite{langlands1971}中系统化了这一构想，并在1970年代与 Jacquet 合作证明了 $\operatorname{GL}_2$ 的第一个非平凡情形\cite{jacquet1970}。此后半个世纪，纲领从数论的一个分支成长为横跨数论、表示论、代数几何乃至数学物理的统一框架——“朗兰兹纲领”一词本身也已成为“数学的宏大统一”的代名词。

\subsection{从数域到几何：纲领的三个层次}

纲领随后沿三个方向铺开，每一层都以函数域/几何情形为试验场：

\begin{itemize}
    \item \textbf{局部朗兰兹}：局部域上 $\operatorname{GL}_n$ 的不可约表示 $\leftrightarrow$ 韦伊群到 $\operatorname{GL}_n(\mathbb{C})$ 的同态（后推广到一般约化群）；
    \item \textbf{整体朗兰兹}：数域或函数域上的尖点自守表示 $\leftrightarrow$ 伽罗瓦/韦伊群表示；
    \item \textbf{几何朗兰兹}：曲线 $\mathbb{P}^1$ 或一般曲线上的主丛范畴 $\operatorname{Bun}_G$ 的 D-模（或拟凝聚层）$\leftrightarrow$ 朗兰兹对偶群的局部系统范畴——一个“范畴化”的对应。
\end{itemize}

这三个层次的历史推进（Drinfeld 1980s、L. Lafforgue 2002、Fargues--Scholze 2021、Gaitsgory--Raskin 2024）将在\cref{sec:results}中详述。
